\documentclass{ximera}
\input{../../../preamble.tex}


\definecolor{fvdnavy}{HTML}{071D43}
\definecolor{fvdcyan}{HTML}{20BFC6}
\definecolor{fvdcyanlight}{HTML}{EFFCFB}
\definecolor{fvdmagenta}{HTML}{F10D69}
\definecolor{fvdmagentalight}{HTML}{FFF1F7}
\definecolor{fvdtext}{HTML}{163044}





%=================================================
% Pedagogische omgevingen
% PDF: tcolorbox
% HTML: learning-box
%=================================================

\newenvironment{voorkennis}
{%
    \ifdefined\HCode
        \HCode{%
<section class="learning-box learning-box--prior">
<div class="learning-box__header">
<span class="learning-box__icon" aria-hidden="true"></span>
<span class="learning-box__title">Voorkennis</span>
</div>
<div class="learning-box__content">
        }%
        \begin{itemize}
    \else
       \begin{tcolorbox}[
    enhanced,
    breakable,
    colback=fvdcyanlight,
    colframe=fvdcyan,
    coltext=fvdtext,
    boxrule=0.5pt,
    leftrule=4pt,
    arc=4mm,
    outer arc=4mm,
    left=7mm,
    right=6mm,
    top=5mm,
    bottom=5mm,
    before skip=6mm,
    after skip=6mm,
    title=Voorkennis,
    coltitle=fvdcyan,
    fonttitle=\bfseries\Large,
    attach title to upper={\par\smallskip},
    boxed title style={
        empty,
        left=0mm,
        right=0mm,
        top=0mm,
        bottom=0mm
    },
    overlay={
        \draw[
            fvdcyan,
            line width=0.5pt
        ]
        ([xshift=42mm,yshift=-13mm]frame.north west)
        --
        ([xshift=-7mm,yshift=-13mm]frame.north east);
    }
]
        \begin{itemize}
    \fi
}
{%
    \end{itemize}
    \ifdefined\HCode
        \HCode{%
</div>
</section>
        }%
    \else
        \end{tcolorbox}
    \fi
}


\newenvironment{leerdoelen}
{%
    \ifdefined\HCode
        \HCode{%
<section class="learning-box learning-box--goals">
<div class="learning-box__header">
<span class="learning-box__icon" aria-hidden="true"></span>
<span class="learning-box__title">Leerdoelen</span>
</div>
<div class="learning-box__content">
        }%
        \begin{itemize}
    \else
\begin{tcolorbox}[
    enhanced,
    breakable,
    colback=fvdmagentalight,
    colframe=fvdmagenta,
    coltext=fvdtext,
    boxrule=0.5pt,
    leftrule=4pt,
    arc=4mm,
    outer arc=4mm,
    left=7mm,
    right=6mm,
    top=5mm,
    bottom=5mm,
    before skip=6mm,
    after skip=6mm,
    title=Leerdoelen,
    coltitle=fvdmagenta,
    fonttitle=\bfseries\Large,
    attach title to upper={\par\smallskip},
    boxed title style={
        empty,
        left=0mm,
        right=0mm,
        top=0mm,
        bottom=0mm
    },
    overlay={
        \draw[
            fvdmagenta,
            line width=0.5pt
        ]
        ([xshift=42mm,yshift=-13mm]frame.north west)
        --
        ([xshift=-7mm,yshift=-13mm]frame.north east);
    }
]
        \begin{itemize}
    \fi
}
{%
    \end{itemize}
    \ifdefined\HCode
        \HCode{%
</div>
</section>
        }%
    \else
        \end{tcolorbox}
    \fi
}


\addPrintStyle{../../..}
\title{Parate kennis --- Je wiskundige gereedschapskist}
\author{Ingmar Herreman}
\begin{abstract}
In dit hoofdstuk brengen we de wiskundige technieken die je uit vorige jaren kent opnieuw op scherp.
Dit is geen hoofdstuk waarin we alles van nul herbeginnen, maar je wiskundige gereedschapskist.
Wat je al vlot beheerst, hoef je niet opnieuw uitgebreid te oefenen. Waar een zelfcheck moeilijk loopt,
oefen je gericht bij. Doorheen de rest van de cursus verwachten we dat je deze technieken zelfstandig
kunt herkennen, inzetten en controleren.
\end{abstract}
\begin{document}
\fvdchapterdata{1}{Verbanden, functies en gedrag}{2}
\maketitle

\begin{voorkennis}
\item De leerstof wiskunde uit de tweede graad.
\item Basisrekenvaardigheid met gehele en rationale getallen.
\item Elementaire algebra en meetkunde.
\end{voorkennis}

\begin{leerdoelen}
\item Je rekent vlot met getallen, breuken, machten en wortels.
\item Je gebruikt wetenschappelijke notatie, verhoudingen en procentuele verandering.
\item Je vereenvoudigt en ontbindt algebraïsche uitdrukkingen.
\item Je lost vergelijkingen en ongelijkheden op en vormt formules om.
\item Je rekent correct met eenheden.
\item Je bepaalt de richtingscoëfficiënt en een vergelijking van een rechte.
\item Je gebruikt coördinaten, Pythagoras en basisgoniometrie.
\item Je kiest zelfstandig geschikt wiskundig gereedschap en controleert je resultaat.
\end{leerdoelen}

\FVDsection{Je wiskundige gereedschapskist}
Dit thema is bewust compact. De oefeningen met \textbf{\ESS} vormen het minimumtraject.
Beheers je een zelfcheck vlot, dan kun je verder; maak je fouten, dan oefen je dat onderdeel gericht bij.

\FVDsection{Rekenen met getallen}
\FVDsubsection{Rekenvolgorde en breuken}
\[
\boxed{\text{haakjes}\to\text{machten en wortels}\to\text{vermenigvuldigen en delen}\to\text{optellen en aftrekken}}
\]
Voor $b,d\neq0$:
\[
\frac ab+\frac cd=\frac{ad+bc}{bd},\qquad
\frac ab\cdot\frac cd=\frac{ac}{bd},\qquad
\frac{\frac ab}{\frac cd}=\frac ab\cdot\frac dc.
\]
\begin{opmerking}
Je mag alleen factoren schrappen, geen termen. Maak bij numerieke berekeningen geregeld eerst een ruwe schatting.
\end{opmerking}
\begin{oefening}[Zelfcheck: rekenen]{\oefB\ESS}
\begin{question}$7-2(3-5)=\answer{11}$.\end{question}
\begin{question}$\frac34+\frac56=\answer{\frac{19}{12}}$.\end{question}
\begin{question}$\frac56:\frac{10}{9}=\answer{\frac34}$.\end{question}
\end{oefening}

\FVDsection{Machten, wortels en wetenschappelijke notatie}
\[
a^ma^n=a^{m+n},\quad \frac{a^m}{a^n}=a^{m-n},\quad
(a^m)^n=a^{mn},\quad a^{-n}=\frac1{a^n},\quad a^0=1.
\]
\[
\sqrt{a^2}=|a|,\qquad \sqrt{72}=6\sqrt2.
\]
Een getal in wetenschappelijke notatie heeft de vorm $a\cdot10^n$ met $1\leq|a|<10$.
\[
0{,}0000032=3{,}2\cdot10^{-6},\qquad
\frac{6\cdot10^8}{2\cdot10^{-3}}=3\cdot10^{11}.
\]
\begin{oefening}[Zelfcheck: machten en wortels]{\oefB\ESS}
\begin{question}$\frac{x^7}{x^3}=\answer{x^4}$.\end{question}
\begin{question}$\sqrt{50}=\answer{5\sqrt2}$.\end{question}
\begin{question}$0{,}00042=\answer{4.2}\cdot10^{\answer{-4}}$.\end{question}
\end{oefening}

\FVDsection{Procenten, verhoudingen en verandering}
Van een oude waarde $O$ naar een nieuwe waarde $N$:
\[
\text{relatieve verandering}=\frac{N-O}{O},\qquad
\text{procentuele verandering}=\frac{N-O}{O}\cdot100\%.
\]
\begin{voorbeeld}
Van $80$ naar $92$ is een stijging van $12$, dus $\frac{12}{80}=0{,}15=15\%$.
\end{voorbeeld}
\begin{oefening}[Zelfcheck: relatieve verandering]{\oefB\ESS}
Een hoeveelheid stijgt van $250$ naar $290$.
\begin{question}De absolute stijging is $\answer{40}$.\end{question}
\begin{question}De procentuele stijging is $\answer{16}\%$.\end{question}
\end{oefening}

\FVDsection{Algebraïsch rekenen}
\[
a(b+c)=ab+ac,
\]
\[
(a+b)^2=a^2+2ab+b^2,\quad
(a-b)^2=a^2-2ab+b^2,\quad
(a-b)(a+b)=a^2-b^2.
\]
Je moet deze regels in beide richtingen herkennen. Ontbinden betekent een som of verschil als product schrijven.
\[
6x^2-15x=3x(2x-5),\qquad x^2-9=(x-3)(x+3).
\]
\begin{oefening}[Zelfcheck: algebra]{\oefB\ESS}
\begin{question}$(2x-3)^2=\answer{4x^2-12x+9}$.\end{question}
\begin{question}$5x^2-20x=\answer{5x(x-4)}$.\end{question}
\begin{question}$x^2-49=\answer{(x-7)(x+7)}$.\end{question}
\end{oefening}

\FVDsection{Vergelijkingen, ongelijkheden en formules}
Een uitdrukking, identiteit en vergelijking zijn verschillende objecten:
\[
3x+2,\qquad (a+b)^2=a^2+2ab+b^2,\qquad 3x+2=11.
\]
\FVDsubsection{Formules omvormen}
In wetenschappen is de onbekende niet altijd $x$. Uit
\[
E_k=\frac12mv^2
\]
volgt voor een niet-negatieve snelheid:
\[
v=\sqrt{\frac{2E_k}{m}}.
\]
\begin{oefening}[Formules omvormen]{\oefB\ESS}
Gegeven $v=\frac{s}{t}$.
\begin{question}Maak $s$ vrij: $s=\answer{vt}$.\end{question}
\begin{question}Maak $t$ vrij: $t=\answer{\frac{s}{v}}$.\end{question}
\end{oefening}

\FVDsubsection{Ongelijkheden en intervalnotatie}
Bij vermenigvuldigen of delen door een negatief getal keert het ongelijkheidsteken om.
\[
-2x<6\iff x>-3\iff x\in]-3,+\infty[.
\]
We gebruiken $[a,b]$, $]a,b[$, $[a,b[$ en $]a,b]$. Oneindig is nooit inbegrepen.
\begin{oefening}[Zelfcheck: ongelijkheden]{\oefB\ESS}
\begin{question}$3x-4\leq8\iff x\leq\answer{4}$.\end{question}
\begin{question}
Welke notatie hoort bij $x>2$?
\begin{multipleChoice}
\choice{$[2,+\infty[$}\choice[correct]{$]2,+\infty[$}\choice{$]-\infty,2[$}
\end{multipleChoice}
\end{question}
\end{oefening}

\FVDsection{Eenheden en grootheden}
\[
72\,\mathrm{km/h}
=72\cdot\frac{1000\,\mathrm m}{3600\,\mathrm s}
=20\,\mathrm{m/s}.
\]
Let op bij oppervlakte en volume:
\[
1\,\mathrm{cm^2}=10^{-4}\,\mathrm{m^2},\qquad
1\,\mathrm{cm^3}=10^{-6}\,\mathrm{m^3}.
\]
Eenheden zijn ook een controle-instrument: een snelheid hoort bijvoorbeeld een eenheid als $\mathrm{m/s}$ te hebben.
\begin{oefening}[Zelfcheck: eenheden]{\oefB\ESS}
\begin{question}$36\,\mathrm{km/h}=\answer{10}\,\mathrm{m/s}$.\end{question}
\begin{question}$1\,\mathrm{cm^2}=10^{\answer{-4}}\,\mathrm{m^2}$.\end{question}
\end{oefening}

\FVDsection{Rechten en coördinaten}
Voor $A(x_1,y_1)$ en $B(x_2,y_2)$ met $x_1\neq x_2$:
\[
m=\frac{y_2-y_1}{x_2-x_1}.
\]
Een niet-verticale rechte kan geschreven worden als
\[
y=mx+q.
\]
Het midden van $[AB]$ is
\[
M\left(\frac{x_1+x_2}{2},\frac{y_1+y_2}{2}\right).
\]
\begin{oefening}[Zelfcheck: rechten]{\oefB\ESS}
Gegeven $A(2,1)$ en $B(6,9)$.
\begin{question}$m=\answer{2}$.\end{question}
\begin{question}$M(\answer{4},\answer{5})$.\end{question}
\end{oefening}

\FVDsection{Pythagoras en basisgoniometrie}
In een rechthoekige driehoek:
\[
a^2+b^2=c^2,
\]
waarbij $c$ de schuine zijde is.
Voor twee punten volgt daaruit:
\[
|AB|=\sqrt{(x_2-x_1)^2+(y_2-y_1)^2}.
\]
Voor een scherpe hoek $\alpha$:
\[
\sin\alpha=\frac{\text{overstaande}}{\text{schuine}},\quad
\cos\alpha=\frac{\text{aanliggende}}{\text{schuine}},\quad
\tan\alpha=\frac{\text{overstaande}}{\text{aanliggende}}.
\]
\begin{oefening}[Zelfcheck: Pythagoras en goniometrie]{\oefB\ESS}
\begin{question}Rechthoekszijden $5$ en $12$: $c=\answer{13}$.\end{question}
\begin{question}Afstand tussen $(0,0)$ en $(3,4)$: $\answer{5}$.\end{question}
\begin{question}Overstaande zijde $3$, aanliggende zijde $4$: $\tan\alpha=\answer{\frac34}$.\end{question}
\end{oefening}

\FVDsection{Kan dit kloppen?}
Controleer na een berekening minstens: teken, ordegrootte, eenheid en voorwaarden.
\begin{oefening}[Fouten opsporen]{\oefU\ESS}
\begin{question}
Een leerling vindt $72\,\mathrm{km/h}=200\,\mathrm{m/s}$.
\begin{multipleChoice}
\choice{Dit kan kloppen.}
\choice[correct]{Dit is onredelijk; de ordegrootte klopt niet.}
\choice{Eenheden mogen niet omgezet worden.}
\end{multipleChoice}
\end{question}
\begin{question}
Een leerling schrijft $\sqrt{50}=25\sqrt2$.
\begin{multipleChoice}
\choice[correct]{$\sqrt{50}$ ligt tussen $7$ en $8$, dus dit kan niet kloppen.}
\choice{Een wortel mag nooit vereenvoudigd worden.}
\end{multipleChoice}
\end{question}
\end{oefening}

\FVDsection{Welk gereedschap heb je nodig?}
Bij een echt probleem staat niet vermeld welke techniek je moet gebruiken.
\[
\boxed{\text{gegeven}\to\text{gevraagd}\to\text{verband}\to\text{gereedschap}\to\text{controle}}
\]
\begin{oefening}[Kies je gereedschap]{\oefU\ESS}
\begin{question}
Een ladder van $5\,\mathrm m$ staat tegen een muur; de voet staat $3\,\mathrm m$ van de muur.
\begin{multipleChoice}
\choice{procentrekening}\choice[correct]{Pythagoras}\choice{wetenschappelijke notatie}
\end{multipleChoice}
\end{question}
\begin{question}
In $E=\frac12mv^2$ zijn $E$ en $m$ gekend en $v$ wordt gevraagd.
\begin{multipleChoice}
\choice{richtingscoëfficiënt}\choice[correct]{formule omvormen}\choice{procentrekening}
\end{multipleChoice}
\end{question}
\end{oefening}

\FVDsection{Synthese: één probleem, meerdere gereedschappen}
\begin{oefening}[Een dronevlucht analyseren]{\oefU\ESS}
Een drone beweegt van $A(0,0)$ naar $B(120,50)$, coördinaten in meter, in $10\,\mathrm s$.
\begin{question}Welke techniek bepaalt de rechte afstand?
\begin{multipleChoice}\choice{procentrekening}\choice[correct]{Pythagoras}\choice{ongelijkheid}\end{multipleChoice}
\end{question}
\begin{question}$|AB|=\answer{130}\,\mathrm m$.\end{question}
\begin{question}Met $v=s/t$: $v=\answer{13}\,\mathrm{m/s}$.\end{question}
\begin{question}$13\,\mathrm{m/s}=\answer{46.8}\,\mathrm{km/h}$.\end{question}
\begin{question}
Een tweede vlucht heeft $15\,\mathrm{m/s}$. De procentuele stijging t.o.v. $13\,\mathrm{m/s}$ is afgerond
$\answer{15.4}\%$.
\end{question}
\end{oefening}

\FVDsection{Essentiële oefeningen}
De oefeningen met \textbf{\ESS} vormen het minimumtraject. Ze controleren rekenen en breuken,
machten en wortels, wetenschappelijke notatie, procentuele verandering, algebra, formules omvormen,
ongelijkheden en intervalnotatie, eenheden, rechten, Pythagoras en goniometrie, controlevaardigheid en
het zelfstandig kiezen en combineren van gereedschappen.

\FVDsection{En nu?}
Je gereedschapskist staat opnieuw klaar. In het volgende thema gebruiken we ze bij meetgegevens,
spreidingsdiagrammen, correlatie en regressie.

\end{document}
